\documentclass[10pt,a4paper]{article}
\usepackage[utf8]{inputenc}
\usepackage[margin=2cm]{geometry}
\usepackage{amsmath}
\usepackage{hyperref}
\pagestyle{empty}
\setlength{\parskip}{0.38em}
\setlength{\parindent}{0pt}

\begin{document}
\noindent
November 15, 2026

\vspace{0.6em}
\noindent
To the Editor-in-Chief and Editorial Board of \textit{Renewable Energy}

\vspace{0.6em}
\noindent
Dear Editor,

\vspace{0.5em}
\noindent
We are pleased to submit our original research article entitled \textit{``Semi-Transparent Organic Agrivoltaics for Cameroonian Perennial Crops: Seasonal Shade Response, Cooperative Techno-Economics, and a Bound on Quantum-Coherence Effects''} for consideration in \textit{Renewable Energy}.

Wavelength-selective semi-transparent organic photovoltaics (ST-OPVs) are widely proposed for tropical agrivoltaics on the premise that transmitted spectra enhance crop productivity relative to opaque crystalline silicon (c-Si) while modulating photosynthetic quantum coherence. Existing evaluations rarely link calibrated transfer-matrix optics, hourly tropical microclimate and stomatal--soil-water physiology, cooperative processing dispatch, replacement-aware life-cycle accounting, and global uncertainty propagation within one self-consistent counterfactual framework. Here, we assess three perennial production chains across western Cameroon's wet--dry highland belt---Arabica coffee (Foumbot, 1,044~m), Protected Geographical Indication (PGI) Penja white pepper (Njombe-Penja, 273~m), and highland tea (Djuttitsa, 1,400~m)---over a full 8,760~h (2023 ERA5) meteorological cycle and a 2014--2023 multi-year reanalysis cross-check.

A four-lever Shapley decomposition of the modeled assimilation shifts attributes a pronounced seasonal regime inversion: dry-season leaf cooling ($\approx 1.5$--$2.7\,^\circ\mathrm{C}$ rainfed, $3.1$--$3.4\,^\circ\mathrm{C}$ under supplemental irrigation $D_5$) boosts dry-season assimilation across all species ($+8.8\%$ to $+30.2\%$ rainfed, up to $+49.7\%$ irrigated), whereas monsoon cloudiness drives net annual assimilation proxies negative in coffee ($-7.08\%$, 2014--2023 decade median $-8.25\%$) and tea ($-3.35\%$, decade median $-3.36\%$) under $40.6\%$ shade, leaving Penja white pepper as the sole crop sustaining net rainfed gains in 2023 ($+7.43\%$ to $+8.16\%$, land equivalent ratio $1.48$--$1.55$; positive in seven of ten reanalysis years, decade median $+2.99\%$, mean $+2.73\%$). At matched $\sim 40\%$ shade, ST-OPV ($D_3$) and elevated c-Si ($D_2$) crop responses differ by at most $1.05$ percentage points, whereas matching annual electricity output ($D_{2,\mathrm{EM}}$, $\mathrm{GCR}=32.4\%$--$33.2\%$) lowers c-Si shade and improves both crop assimilation and net present value (NPV), confirming within the model that shade fraction, service basis, and seasonality---not spectral selectivity---govern agronomy. Multi-Hamiltonian continuous-wave steady-state bounds (cross-checked against MesoHOPS~1.8.0 pulsed propagation and 4-site exact HEOM) constrain excitonic efficiency shifts below $3.37\times 10^{-5}$ ($2.10\times 10^{-3}$ HEOM envelope), resulting in negligible Sobol sensitivity ($S_{Ti} \le 2.59\times 10^{-7}$).

Across $N=10,000$ Saltelli quasi-Monte Carlo base draws, a dual variance decomposition shows that grid export tariffs and cost of capital govern standalone ST-OPV bankability, whereas hardware costs ($C_{\mathrm{mod}}^{\mathrm{OPV}}, C_{\mathrm{BOS}}^{\mathrm{c\text{-}Si}}$) and encapsulation lifetime ($T_{80}$) govern competition against elevated c-Si, which achieves higher 20-year NPV in $77.0\%$--$89.0\%$ of draws ($75.3\%$--$87.7\%$ at zero carbon). Life-cycle carbon parity requires $T_{80,\mathrm{LCA}}^* \ge 7.40$--$8.12$~yr (slashing intensity to $7.8$--$12.5~\mathrm{gCO_2e/kWh}$ at 25~yr), whereas economic NPV parity emerges only in Penja pepper at $\$15/\mathrm{m}^2$ with $T_{80,\mathrm{NPV}}^* \ge 16.33$~yr; overnight tea withering storage lifts solar fraction to $99.9\%$ and sustains positive off-grid NPV ($+\$319,794/\mathrm{ha}$), while load-matched off-grid downsizing ($0.07$--$0.39$~ha) recovers positive returns across all three crops ($+\$160.3\mathrm{k}$--$+\$622.8\mathrm{k}/\mathrm{ha}$ over open sun).

The manuscript runs to 20 pages---4,705 words of clean main text (5,194 including captions) with 50 references---and is accompanied by a 13-page Supporting Material, a Graphical Abstract, and Research Highlights.

This manuscript is original, has not been published previously, and is not under consideration elsewhere. All authors approve this submission and declare no competing financial interests.

\vspace{0.6em}
\noindent
Sincerely,

\vspace{0.5em}
\noindent
Serge Guy Nana Engo (Corresponding author, on behalf of all co-authors)\\
Department of Physics, Faculty of Science, University of Yaound\'e~I, P.O. Box 812, Yaound\'e, Cameroon\\
E-mail: serge.nana-engo@facsciences-uy1.cm
\end{document}
